% GENERATED FILE. Edit canonical lessons, then run npm run generate:books.
% canonical-source-hash: fnv1a64:783d75266713a000
% canonical-lessons: RU-C280-soobshchat, RU-C280-otritsat, RU-C280-publikovat, RU-C280-podtverzhdat, RU-C280-oprovergat

\chapter{Reporting and denying}
\label{ch:ru-reporting-and-denying}
\hlchaptermodality{280}

\begin{chapteropening}
\textbf{By the end of this chapter:} \emph{I can say that someone reported, denied, published, confirmed or refuted something.}

Complete the last lesson of chapter 280: I can say that someone reported, denied, published, confirmed or refuted something.
\end{chapteropening}

\section[\texorpdfstring{\ru{сообщать} (soobshchát)}{сообщать (soobshchát)}]{\ru{сообщать} (soobshchát) — to report, to inform}

\label{lesson:RU-C280-soobshchat}



\begin{quote}
Before the new one: say the Russian for an article, then the Russian for a sensation.
\end{quote}

\subsection*{You'll want to know: \ru{сообщать}}
\textbf{\ru{сообщать}} (\emph{soobshchát}) — \textquotedblleft{}to report, to inform\textquotedblright{}.

\begin{quote}
\ru{Радио} \ru{сообщает} \ru{о} \ru{сильном} \ru{ветре}.
\end{quote}

\begin{quote}
\emph{Rádio soobshcháyet o sílnom vétre.}
\end{quote}

\begin{quote}
The radio is reporting a strong wind.
\end{quote}

Imperfective; perfective \emph{\ru{сообщить}} (\emph{soobshchít}). The person is in the dative, the news in the accusative or with \emph{\ru{о}}: \emph{\ru{сообщить} \ru{маме} \ru{новость}}.

\begin{cousinweb}
\emph{\ru{сообщать}} is built on \emph{\ru{со}-} and \emph{\ru{общий}} \textquotedblleft{}common\textquotedblright{}, so \textquotedblleft{}to make shared\textquotedblright{}; \emph{\ru{общий}} is itself a Church Slavonic form (EtymDB: from Church Slavonic \emph{obĭščĭ}, Proto-Slavic \emph{obĭťĭ} \textquotedblleft{}common\textquotedblright{}), and EtymDB lists no parent for the verb.
\end{cousinweb}

\begin{grammarlens}[title={What you've built}]
One more word for telling, explaining and describing at B1.
\end{grammarlens}

\subsection*{Your turn}
\begin{itemize}
\raggedright
  \item \emph{Say it:} \emph{soobshchát}
  \item \emph{Say it:} \emph{soobshchát}, once more
  \item \emph{Say it:} a sentence of your own with \emph{soobshchát}
  \item \emph{Recall:} say the Russian for an article, then the Russian for a sensation, then say \emph{soobshchát} again
\end{itemize}

\begin{tcolorbox}[breakable,colback=teal!4,colframe=teal!35!black,title={Before you move on}]
What does \ru{сообщать} mean? (\textquotedblleft{}To report, to inform\textquotedblright{}.) Say it once more.
\end{tcolorbox}
\section[\texorpdfstring{\ru{отрицать} (otritsát)}{отрицать (otritsát)}]{\ru{отрицать} (otritsát) — to deny}

\label{lesson:RU-C280-otritsat}



\begin{quote}
Before the new one: say the Russian for a sensation, then the Russian for to report.
\end{quote}

\subsection*{You'll want to know: \ru{отрицать}}
\textbf{\ru{отрицать}} (\emph{otritsát}) — \textquotedblleft{}to deny\textquotedblright{}.

\begin{quote}
\ru{Он} \ru{отрицает}, \ru{что} \ru{был} \ru{там}.
\end{quote}

\begin{quote}
\emph{On otritsáyet, shto byl tam.}
\end{quote}

\begin{quote}
He denies that he was there.
\end{quote}

Imperfective, with no everyday perfective. It takes the accusative: \emph{\ru{отрицать} \ru{всё}} (\emph{otritsát vsyo}) — to deny everything.

\begin{cousinweb}
\emph{\ru{отрицать}} is usually described as a Church Slavonic verb, traced to \emph{\ru{от}-} \textquotedblleft{}away, against\textquotedblright{} and the root of \emph{\ru{рещи}} \textquotedblleft{}to say\textquotedblright{}, so \textquotedblleft{}to speak against\textquotedblright{}; EtymDB lists no parent.
\end{cousinweb}

\begin{grammarlens}[title={What you've built}]
One more word for telling, explaining and describing at B1.
\end{grammarlens}

\subsection*{Your turn}
\begin{itemize}
\raggedright
  \item \emph{Say it:} \emph{otritsát}
  \item \emph{Say it:} \emph{otritsát}, once more
  \item \emph{Say it:} a sentence of your own with \emph{otritsát}
  \item \emph{Recall:} say the Russian for a sensation, then the Russian for to report, then say \emph{otritsát} again
\end{itemize}

\begin{tcolorbox}[breakable,colback=teal!4,colframe=teal!35!black,title={Before you move on}]
What does \ru{отрицать} mean? (\textquotedblleft{}To deny\textquotedblright{}.) Say it once more.
\end{tcolorbox}
\section[\texorpdfstring{\ru{публиковать} (publikovát)}{публиковать (publikovát)}]{\ru{публиковать} (publikovát) — to publish}

\label{lesson:RU-C280-publikovat}



\begin{quote}
Before the new one: say the Russian for to report, then the Russian for to deny.
\end{quote}

\subsection*{You'll want to know: \ru{публиковать}}
\textbf{\ru{публиковать}} (\emph{publikovát}) — \textquotedblleft{}to publish\textquotedblright{}.

\begin{quote}
\ru{Газета} \ru{публикует} \ru{эти} \ru{письма}.
\end{quote}

\begin{quote}
\emph{Gazéta publikúyet éti písma.}
\end{quote}

\begin{quote}
The newspaper publishes these letters.
\end{quote}

Imperfective; the perfective is \emph{\ru{опубликовать}} (\emph{opublikovát}). Present \emph{\ru{публикую}}, \emph{\ru{публикуешь}}.

\begin{cousinweb}
\emph{\ru{публиковать}} goes back to Latin \emph{pūblicāre} \textquotedblleft{}to make public\textquotedblright{}, from \emph{pūblicus} \textquotedblleft{}of the people\textquotedblright{} (Lewis \& Short), with the Russian verb ending \emph{-\ru{овать}}; EtymDB lists the Russian word with no parent.
\end{cousinweb}

\begin{grammarlens}[title={What you've built}]
One more word for telling, explaining and describing at B1.
\end{grammarlens}

\subsection*{Your turn}
\begin{itemize}
\raggedright
  \item \emph{Say it:} \emph{publikovát}
  \item \emph{Say it:} \emph{publikovát}, once more
  \item \emph{Say it:} a sentence of your own with \emph{publikovát}
  \item \emph{Recall:} say the Russian for to report, then the Russian for to deny, then say \emph{publikovát} again
\end{itemize}

\begin{tcolorbox}[breakable,colback=teal!4,colframe=teal!35!black,title={Before you move on}]
What does \ru{публиковать} mean? (\textquotedblleft{}To publish\textquotedblright{}.) Say it once more.
\end{tcolorbox}
\section[\texorpdfstring{\ru{подтверждать} (podtverzhdát)}{подтверждать (podtverzhdát)}]{\ru{подтверждать} (podtverzhdát) — to confirm}

\label{lesson:RU-C280-podtverzhdat}



\begin{quote}
Before the new one: say the Russian for to deny, then the Russian for to publish.
\end{quote}

\subsection*{You'll want to know: \ru{подтверждать}}
\textbf{\ru{подтверждать}} (\emph{podtverzhdát}) — \textquotedblleft{}to confirm\textquotedblright{}.

\begin{quote}
\ru{Врач} \ru{подтвердил}, \ru{что} \ru{всё} \ru{в} \ru{порядке}.
\end{quote}

\begin{quote}
\emph{Vrach podtverdíl, shto vsyo v poryádke.}
\end{quote}

\begin{quote}
The doctor confirmed that everything was all right.
\end{quote}

Imperfective; perfective \emph{\ru{подтвердить}} (\emph{podtverdít}), \emph{\ru{я} \ru{подтвержу}}.

\begin{cousinweb}
\emph{\ru{подтверждать}} is built on \emph{\ru{под}-} and the root of \emph{\ru{твёрдый}} \textquotedblleft{}firm\textquotedblright{}, Proto-Slavic \emph{tvĭrdŭ}, in its Church Slavonic shape \emph{-\ru{твержд}-}: to make firm. EtymDB lists no parent for the verb.
\end{cousinweb}

\begin{grammarlens}[title={What you've built}]
One more word for telling, explaining and describing at B1.
\end{grammarlens}

\subsection*{Your turn}
\begin{itemize}
\raggedright
  \item \emph{Say it:} \emph{podtverzhdát}
  \item \emph{Say it:} \emph{podtverzhdát}, once more
  \item \emph{Say it:} a sentence of your own with \emph{podtverzhdát}
  \item \emph{Recall:} say the Russian for to deny, then the Russian for to publish, then say \emph{podtverzhdát} again
\end{itemize}

\begin{tcolorbox}[breakable,colback=teal!4,colframe=teal!35!black,title={Before you move on}]
What does \ru{подтверждать} mean? (\textquotedblleft{}To confirm\textquotedblright{}.) Say it once more.
\end{tcolorbox}
\section[\texorpdfstring{\ru{опровергать} (oprovergát)}{опровергать (oprovergát)}]{\ru{опровергать} (oprovergát) — to refute, to deny officially}

\label{lesson:RU-C280-oprovergat}



\begin{quote}
Before the new one: say the Russian for to publish, then the Russian for to confirm.
\end{quote}

\subsection*{You'll want to know: \ru{опровергать}}
\textbf{\ru{опровергать}} (\emph{oprovergát}) — \textquotedblleft{}to refute, to deny officially\textquotedblright{}.

\begin{quote}
\ru{Компания} \ru{опровергает} \ru{эту} \ru{информацию}.
\end{quote}

\begin{quote}
\emph{Kompániya oprovergáyet étu informátsiyu.}
\end{quote}

\begin{quote}
The company refutes this information.
\end{quote}

Imperfective; perfective \emph{\ru{опровергнуть}} (\emph{oprovérgnut}). Stronger than \emph{\ru{отрицать}}: to show something is false.

\begin{cousinweb}
\emph{\ru{опровергать}} is usually described as a Church Slavonic verb, built on \emph{\ru{о}-}, \emph{\ru{про}-} and \emph{\ru{вергнуть}} \textquotedblleft{}to throw\textquotedblright{}, so \textquotedblleft{}to throw down\textquotedblright{}; EtymDB lists no parent.
\end{cousinweb}

\begin{grammarlens}[title={What you've built}]
One more word for telling, explaining and describing at B1.
\end{grammarlens}

\subsection*{Your turn}
\begin{itemize}
\raggedright
  \item \emph{Say it:} \emph{oprovergát}
  \item \emph{Say it:} \emph{oprovergát}, once more
  \item \emph{Say it:} a sentence of your own with \emph{oprovergát}
  \item \emph{Recall:} say the Russian for to publish, then the Russian for to confirm, then say \emph{oprovergát} again
\end{itemize}

\begin{tcolorbox}[breakable,colback=teal!4,colframe=teal!35!black,title={Before you move on}]
What does \ru{опровергать} mean? (\textquotedblleft{}To refute, to deny officially\textquotedblright{}.) Say it once more.
\end{tcolorbox}
